% ========== KERNEL BOOTSTRAP PROCESS ==========
% Initialization phases, dependency resolution, and failure handling
% Source: Symphony/Content/Microkernel Architecture Guide, The Kernels

\section{Kernel Bootstrap Process}
\label{sec:kernel-bootstrap}

\lettrine{T}{he} bootstrap process represents one of the most critical phases in Symphony's lifecycle. During these initial moments, the system must carefully orchestrate the startup of dozens of components while handling dependencies, managing failures, and ensuring that the system reaches a stable, operational state. Symphony's bootstrap process is designed for both reliability and speed, typically completing full system initialization in under two seconds.

\subsection{Initialization Phases}
\label{subsec:initialization-phases}

Symphony's bootstrap process follows a carefully orchestrated sequence of five distinct phases, each with specific responsibilities and success criteria.

\subsubsection{Phase 1: Core Kernel Initialization}

The bootstrap process begins with the most fundamental system components that form the foundation for all subsequent operations.

\textbf{Conductor Core Startup}

The Python-based Conductor Core initializes first, establishing the basic runtime environment:

\begin{compactlist}
    \item \textbf{Memory Management}: Initialize heap and garbage collection settings
    \item \textbf{Logging System}: Establish structured logging with appropriate levels
    \item \textbf{Configuration Loading}: Read system configuration and validate settings
    \item \textbf{Security Context}: Establish security policies and permission frameworks
    \item \textbf{IPC Foundation}: Initialize the basic IPC infrastructure
\end{compactlist}

\textbf{Critical System Validation}

Before proceeding to the next phase, the system performs essential validation checks:

\begin{table}[h]
\centering
\caption{Phase 1 Validation Checkpoints}
\label{tab:phase1-validation}
\begin{tabular}{@{}lll@{}}
\toprule
\textbf{Validation Check} & \textbf{Success Criteria} & \textbf{Failure Action} \\
\midrule
Memory availability & >512MB free memory & Abort with resource error \\
Configuration validity & All required settings present & Load default configuration \\
Security permissions & Required privileges available & Prompt for elevation \\
IPC socket creation & Unix socket successfully bound & Retry with alternate path \\
Logging system & Log file writable & Fall back to console logging \\
\bottomrule
\end{tabular}
\end{table}

\subsubsection{Phase 2: The Pit Infrastructure}

Once the Conductor Core is stable, the system initializes The Pit extensions that provide high-performance infrastructure services.

\textbf{Sequential Pit Extension Loading}

The five Pit extensions are loaded in a specific order based on their dependencies:

\begin{enumerate}
    \item \textbf{Pool Manager}: Must be first as other extensions may need resource allocation
    \item \textbf{Artifact Store}: Required for storing initialization state and logs
    \item \textbf{DAG Tracker}: Needed for tracking the bootstrap process itself
    \item \textbf{Arbitration Engine}: Required for resolving resource conflicts during startup
    \item \textbf{Stale Manager}: Last, as it needs visibility into all other extension states
\end{enumerate}

\textbf{Shared Memory Initialization}

The Pit extensions establish shared memory regions for high-performance communication:

\begin{infobox}[title=Shared Memory Layout]
The Pit uses a carefully designed shared memory layout with atomic counters, read-optimized structures, and cache-line alignment for optimal performance. Total shared memory allocation: 64MB baseline, expandable to 1GB under load.
\end{infobox}

\textbf{Health Check Establishment}

Each Pit extension implements a health check interface that reports:
\begin{compactlist}
    \item \textbf{Initialization Status}: Whether the extension is fully operational
    \item \textbf{Resource Usage}: Current memory and CPU utilization
    \item \textbf{Performance Metrics}: Response times and throughput capabilities
    \item \textbf{Error Conditions}: Any warnings or degraded functionality
\end{compactlist}

\subsubsection{Phase 3: IPC Bus Activation}

With The Pit operational, the system activates the full IPC infrastructure for out-of-process communication.

\textbf{Transport Layer Initialization}

The IPC bus initializes multiple transport mechanisms based on the target platform:

\begin{table}[h]
\centering
\caption{IPC Transport Initialization by Platform}
\label{tab:ipc-transport-init}
\begin{tabular}{@{}lll@{}}
\toprule
\textbf{Platform} & \textbf{Primary Transport} & \textbf{Fallback Transport} \\
\midrule
Linux & Unix Domain Sockets & Named Pipes \\
macOS & Unix Domain Sockets & Mach Ports \\
Windows & Named Pipes & TCP Loopback \\
\bottomrule
\end{tabular}
\end{table}

\textbf{Message Routing Engine}

The routing engine establishes the message dispatch system that will handle all inter-extension communication:

\begin{compactlist}
    \item \textbf{Routing Tables}: Initialize extension address mappings
    \item \textbf{Load Balancing}: Set up round-robin and weighted routing
    \item \textbf{Circuit Breakers}: Configure failure detection and isolation
    \item \textbf{Message Queues}: Establish buffering for high-throughput scenarios
\end{compactlist}

\subsubsection{Phase 4: Essential Extensions}

The system now loads essential out-of-process extensions that provide core functionality.

\textbf{Core Extension Categories}

Essential extensions are loaded in priority order:

\begin{enumerate}
    \item \textbf{File System Extensions}: File explorer, project management
    \item \textbf{Editor Extensions}: Text editing, syntax highlighting
    \item \textbf{Terminal Extensions}: Integrated terminal, command execution
    \item \textbf{Settings Extensions}: Configuration management, preferences
\end{enumerate}

\textbf{Extension Process Management}

Each extension is launched as an isolated process with:

\begin{compactlist}
    \item \textbf{Resource Limits}: Memory and CPU quotas based on extension manifest
    \item \textbf{Security Context}: Minimal privileges required for functionality
    \item \textbf{Health Monitoring}: Continuous process health and performance tracking
    \item \textbf{Restart Policies}: Automatic restart configuration for different failure types
\end{compactlist}

\subsubsection{Phase 5: User Interface Activation}

The final phase brings up the user interface and makes the system ready for user interaction.

\textbf{React Application Startup}

The React-based user interface initializes with:

\begin{compactlist}
    \item \textbf{Component Registry}: Registration of all available UI components
    \item \textbf{Theme System}: Loading of user preferences and visual themes
    \item \textbf{Event Handlers}: Establishment of user input processing
    \item \textbf{WebSocket Connections}: Real-time communication with backend extensions
\end{compactlist}

\textbf{Initial Workspace Loading}

The system loads the user's workspace and project state:

\begin{compactlist}
    \item \textbf{Project Discovery}: Scan for and load recent projects
    \item \textbf{Session Restoration}: Restore previous editor state and open files
    \item \textbf{Extension Preferences}: Load user-specific extension configurations
    \item \textbf{Layout Restoration}: Restore window layout and panel arrangements
\end{compactlist}

\subsection{Dependency Resolution}
\label{subsec:dependency-resolution}

Symphony's dependency resolution system ensures that extensions are loaded in the correct order and that all dependencies are satisfied before an extension becomes operational.

\subsubsection{Dependency Graph Construction}

The system builds a directed acyclic graph (DAG) of extension dependencies during the bootstrap process:

\begin{figure}[h]
\centering
\begin{tikzpicture}[node distance=2cm, auto]
    % Nodes
    \node[draw, circle] (conductor) {Conductor};
    \node[draw, circle, below left=of conductor] (pool) {Pool Mgr};
    \node[draw, circle, below=of conductor] (artifact) {Artifact Store};
    \node[draw, circle, below right=of conductor] (dag) {DAG Tracker};
    \node[draw, circle, below=of pool] (ipc) {IPC Bus};
    \node[draw, circle, below=of artifact] (file) {File Explorer};
    \node[draw, circle, below=of dag] (editor) {Editor};
    \node[draw, circle, below=of ipc] (ui) {UI System};
    
    % Edges
    \draw[->] (conductor) -> (pool);
    \draw[->] (conductor) -> (artifact);
    \draw[->] (conductor) -> (dag);
    \draw[->] (pool) -> (ipc);
    \draw[->] (artifact) -> (ipc);
    \draw[->] (dag) -> (ipc);
    \draw[->] (ipc) -> (file);
    \draw[->] (ipc) -> (editor);
    \draw[->] (file) -> (ui);
    \draw[->] (editor) -> (ui);
\end{tikzpicture}
\caption{Extension Dependency Graph}
\label{fig:dependency-graph}
\end{figure}

\subsubsection{Topological Sorting}

The bootstrap system uses topological sorting to determine the optimal loading order:

\begin{alertbox}
\textbf{Circular Dependency Detection}: The system automatically detects circular dependencies and reports them as configuration errors. Extensions with circular dependencies cannot be loaded and must be reconfigured.
\end{alertbox}

\textbf{Dependency Resolution Algorithm}

The resolution process follows these steps:

\begin{enumerate}
    \item \textbf{Graph Construction}: Build dependency graph from extension manifests
    \item \textbf{Cycle Detection}: Use depth-first search to identify circular dependencies
    \item \textbf{Topological Sort}: Generate loading order that satisfies all dependencies
    \item \textbf{Parallel Optimization}: Identify extensions that can load concurrently
    \item \textbf{Critical Path Analysis}: Optimize loading sequence for minimum total time
\end{enumerate}

\subsubsection{Dynamic Dependency Handling}

Some dependencies are discovered dynamically during runtime:

\textbf{Optional Dependencies}

Extensions can declare optional dependencies that enhance functionality but are not required:

\begin{compactlist}
    \item \textbf{Graceful Degradation}: Extension operates with reduced functionality if optional dependencies are unavailable
    \item \textbf{Runtime Discovery}: Extensions can discover and utilize optional dependencies that become available later
    \item \textbf{Feature Flags}: Functionality is enabled/disabled based on dependency availability
\end{compactlist}

\textbf{Version Compatibility}

The system handles version compatibility through semantic versioning:

\begin{table}[h]
\centering
\caption{Dependency Version Resolution}
\label{tab:version-resolution}
\begin{tabular}{@{}lll@{}}
\toprule
\textbf{Version Constraint} & \textbf{Resolution Strategy} & \textbf{Fallback Action} \\
\midrule
Exact match (=1.2.3) & Load exact version only & Fail if unavailable \\
Compatible (^1.2.0) & Load highest compatible version & Load minimum compatible \\
Minimum (>=1.0.0) & Load highest available version & Load minimum version \\
Range (1.0.0-2.0.0) & Load highest in range & Load lowest in range \\
\bottomrule
\end{tabular}
\end{table}

\subsection{Failure Handling and Rollback}
\label{subsec:failure-handling-rollback}

Symphony's bootstrap process is designed to handle failures gracefully and provide clear recovery paths when components fail to initialize properly.

\subsubsection{Failure Classification}

The system classifies bootstrap failures into different categories that require different response strategies:

\textbf{Critical Failures}

Failures that prevent the system from reaching a usable state:

\begin{compactlist}
    \item \textbf{Conductor Core Failure}: Complete system abort with diagnostic information
    \item \textbf{Memory Exhaustion}: Immediate shutdown with resource usage report
    \item \textbf{Security Violation}: Abort with security audit log
    \item \textbf{Configuration Corruption}: Attempt recovery with default configuration
\end{compactlist}

\textbf{Recoverable Failures}

Failures that can be worked around while maintaining system functionality:

\begin{compactlist}
    \item \textbf{Extension Load Failure}: Continue without the failed extension
    \item \textbf{Network Unavailability}: Operate in offline mode
    \item \textbf{Permission Denied}: Request elevated privileges or operate with reduced functionality
    \item \textbf{Resource Contention}: Retry with exponential backoff
\end{compactlist}

\textbf{Degraded Operation}

Conditions that reduce functionality but allow continued operation:

\begin{compactlist}
    \item \textbf{Performance Degradation}: Continue with performance warnings
    \item \textbf{Feature Unavailability}: Disable affected features and notify user
    \item \textbf{Reduced Capacity}: Operate with lower resource limits
    \item \textbf{Compatibility Issues}: Use fallback implementations
\end{compactlist}

\subsubsection{Rollback Mechanisms}

When failures occur, Symphony implements several rollback mechanisms to restore system stability:

\textbf{Checkpoint-Based Rollback}

The system creates checkpoints at each bootstrap phase:

\begin{table}[h]
\centering
\caption{Bootstrap Checkpoints and Rollback Actions}
\label{tab:bootstrap-checkpoints}
\begin{tabular}{@{}lll@{}}
\toprule
\textbf{Checkpoint} & \textbf{System State} & \textbf{Rollback Action} \\
\midrule
Phase 1 Complete & Conductor Core operational & Restart with safe configuration \\
Phase 2 Complete & Pit extensions loaded & Restart Pit with minimal configuration \\
Phase 3 Complete & IPC system active & Reset IPC and reload extensions \\
Phase 4 Complete & Essential extensions loaded & Restart failed extensions only \\
Phase 5 Complete & UI system active & Refresh UI without full restart \\
\bottomrule
\end{tabular}
\end{table}

\textbf{Incremental Rollback}

Rather than rolling back the entire system, Symphony attempts incremental rollback:

\begin{enumerate}
    \item \textbf{Component Isolation}: Identify the specific failed component
    \item \textbf{Dependency Analysis}: Determine which components depend on the failed component
    \item \textbf{Selective Restart}: Restart only the failed component and its dependents
    \item \textbf{State Preservation}: Maintain state for unaffected components
    \item \textbf{Progressive Recovery}: Gradually restore full functionality
\end{enumerate}

\subsubsection{Recovery Strategies}

Symphony implements multiple recovery strategies based on the type and severity of failures:

\textbf{Automatic Recovery}

For common, transient failures:

\begin{compactlist}
    \item \textbf{Retry with Backoff}: Exponential backoff for resource contention
    \item \textbf{Alternative Paths}: Try different initialization sequences
    \item \textbf{Reduced Configuration}: Start with minimal settings and expand gradually
    \item \textbf{Safe Mode}: Boot with only essential components
\end{compactlist}

\textbf{User-Assisted Recovery}

For failures requiring user intervention:

\begin{compactlist}
    \item \textbf{Configuration Repair}: Guide user through configuration fixes
    \item \textbf{Permission Elevation}: Request additional privileges when needed
    \item \textbf{Extension Selection}: Allow user to disable problematic extensions
    \item \textbf{Diagnostic Mode}: Provide detailed diagnostic information for troubleshooting
\end{compactlist}

\textbf{Emergency Recovery}

For severe failures that prevent normal operation:

\begin{compactlist}
    \item \textbf{Factory Reset}: Restore system to default configuration
    \item \textbf{Safe Mode Boot}: Start with absolute minimum functionality
    \item \textbf{Diagnostic Collection}: Gather system information for support
    \item \textbf{Backup Restoration}: Restore from previous working configuration
\end{compactlist}

\subsection{Health Checking}
\label{subsec:health-checking}

Symphony implements comprehensive health checking throughout the bootstrap process and during normal operation to ensure system reliability and performance.

\subsubsection{Multi-Level Health Monitoring}

The health checking system operates at multiple levels of granularity:

\textbf{System-Level Health}

Overall system health indicators:

\begin{compactlist}
    \item \textbf{Resource Utilization}: CPU, memory, and disk usage monitoring
    \item \textbf{Performance Metrics}: Response times and throughput measurements
    \item \textbf{Error Rates}: Frequency and severity of system errors
    \item \textbf{Availability}: Uptime and service availability statistics
\end{compactlist}

\textbf{Component-Level Health}

Individual component health monitoring:

\begin{compactlist}
    \item \textbf{Process Health}: Process status, memory leaks, and crash detection
    \item \textbf{Functional Health}: Component-specific functionality verification
    \item \textbf{Performance Health}: Component response times and resource usage
    \item \textbf{Integration Health}: Inter-component communication status
\end{compactlist}

\textbf{Extension-Level Health}

Detailed extension health tracking:

\begin{compactlist}
    \item \textbf{Lifecycle Health}: Extension startup, operation, and shutdown status
    \item \textbf{API Health}: Extension API response times and error rates
    \item \textbf{Resource Health}: Extension resource consumption and limits
    \item \textbf{User Health}: User-reported issues and satisfaction metrics
\end{compactlist}

\subsubsection{Health Check Implementation}

Each component implements standardized health check interfaces:

\textbf{Health Check Protocol}

All components respond to health check requests with structured information:

\begin{table}[h]
\centering
\caption{Health Check Response Format}
\label{tab:health-check-format}
\begin{tabular}{@{}lll@{}}
\toprule
\textbf{Field} & \textbf{Type} & \textbf{Description} \\
\midrule
status & enum & HEALTHY, DEGRADED, UNHEALTHY, UNKNOWN \\
timestamp & datetime & When the health check was performed \\
response\_time & duration & Time taken to perform the health check \\
details & object & Component-specific health information \\
metrics & object & Performance and resource usage metrics \\
errors & array & Current error conditions and warnings \\
\bottomrule
\end{tabular}
\end{table}

\textbf{Proactive Health Monitoring}

The system performs continuous health monitoring:

\begin{compactlist}
    \item \textbf{Periodic Checks}: Regular health verification on configurable intervals
    \item \textbf{Threshold Monitoring}: Automatic alerts when metrics exceed thresholds
    \item \textbf{Trend Analysis}: Detection of gradual performance degradation
    \item \textbf{Predictive Alerts}: Early warning of potential failures
\end{compactlist}

\subsubsection{Health-Based Decision Making}

Health information drives automatic system optimization and recovery decisions:

\textbf{Load Balancing}

Health metrics influence request routing and resource allocation:

\begin{compactlist}
    \item \textbf{Healthy Components}: Receive normal traffic allocation
    \item \textbf{Degraded Components}: Receive reduced traffic with monitoring
    \item \textbf{Unhealthy Components}: Removed from rotation until recovery
    \item \textbf{Unknown Components}: Treated as degraded until status clarified
\end{compactlist}

\textbf{Automatic Recovery}

Health status triggers automatic recovery actions:

\begin{compactlist}
    \item \textbf{Restart Policies}: Automatic restart based on health degradation patterns
    \item \textbf{Resource Adjustment}: Dynamic resource allocation based on health metrics
    \item \textbf{Failover Activation}: Automatic failover to backup components
    \item \textbf{Alert Generation}: Notification of administrators for manual intervention
\end{compactlist}

Through this comprehensive bootstrap process, Symphony ensures reliable system initialization while providing robust failure handling and recovery mechanisms. The health checking system maintains continuous visibility into system status, enabling proactive maintenance and optimization.